\subsection{A second domain: the iterated local search for the travelling salesman problem}
\label{sec:tsp}
\paragraph{Design} Every host so far is a scheduling problem. To test whether the findings belong to scheduling, we applied the protocol unchanged to an iterated local search for the Euclidean TSP: a tour improved by 2-opt and Or-opt with neighbour lists, a perturbation chosen at each iteration by the selector (a local double-bridge, a set of short random reversals, or the shuffle of a segment; five sizes of each) and acceptance of a candidate that is not worse. The action spaces are T5 (one perturbation, five sizes) and T15 (three perturbations $\times$ five sizes). Instances are uniform or clustered points with $n\in\{200,500,1000,2000\}$: eight classes, $12$ tuning and $50$ test instances per class, $T=1000$, three seeds; the learners keep the hyperparameters frozen from the other hosts, and the reward is the improvement of the tour. The validated action of each class is the best fixed action on the tuning instances; all $15$ fixed actions were also run on the test instances. The ten predictions $P_1$--$P_{10}$ and the hashes were recorded before the test run.

\paragraph{Results} The picture of the scheduling hosts repeats (Table~\ref{tab:tsp}). The validated action gains over random choice beyond the margin in all $16$ cells ($0.05$ to $0.90$ points of the tour length; its mean gap to the best tour found by any method is $0.26\%$ in T15, against $0.77\%$ for random choice), and none of the six online learners beats it beyond the margin in any of the $96$ learner--cell pairs; it is better than the standard Q-Learning, the permuted Q-Learning, the policy gradient, the ALNS roulette and the optimistic Q-Learning in all $16$ cells, than UCB1 in $15$ and than Thompson sampling in $12$. The standard Q-Learning loses to random choice in every cell ($0.13$ to $0.63$ points) and the lock-in of Section~\ref{sec:order} reproduces: it puts $61\%$ of its choices in T15 ($65\%$ in T5) on the action listed first, evaluates $0.64$ times as many moves as random choice, and the permuted version (which evaluates $1.11$ times as many) gains in no cell and differs from it beyond the margin in $15$ of $16$ cells. The learners that are sound by design do gain in some cells: Thompson sampling in $7$ of $16$ (up to $0.38$ points), the optimistic Q-Learning in $4$, UCB1 and the policy gradient in one each; but the best online learner of a cell recovers at most half of the validated gain in $12$ of the $16$ cells. All ten predictions held.

\paragraph{Published instances} We then ran the same host on $37$ instances of the TSPLIB~\cite{Reinelt1991} (EUC\_2D, $200\le n\le2103$; Table~\ref{tab:tsplib}; the optimal tour lengths are those of the TSPLIB file of solutions). The static actions were chosen only on the synthetic tuning instances (the action of the synthetic size nearest to the instance, and one action for all), the learners kept their hyperparameters, and nine predictions $L_1$--$L_9$ were recorded beforehand; all nine held. The static action chosen on synthetic instances beats random choice beyond the margin in both spaces ($0.53$ and $0.56$ points of the tour length; one action for all gains $0.53$ and $0.68$), no online learner beats it, and the best online learner, Thompson sampling, recovers $23\%$ and $43\%$ of its gain and is the only one that beats random choice beyond the margin ($+0.24$ in T15). The standard Q-Learning loses ($-0.61$ and $-0.34$) and puts $63\%$ of its choices in T15 on the action listed first; the permuted version loses less ($-0.15$ and $-0.19$) and differs from it beyond the margin in both spaces. The host is close to optimal in absolute terms: the static action ends on average $1.21\%$ above the optimal tour in T15 (random choice: $1.79\%$), $0.76\%$ and $0.72\%$ for the small and medium instances and $2.17\%$ for the large ones, where the static action gains most ($1.19$ points; none in the small ones, where there is no headroom).

\paragraph{The hierarchy of knowledge is not universal} In the TSP the finest knowledge adds nothing: one action chosen for all instances ($G_0$) gains $0.54$ points in T15, more than the class-validated action ($0.50$, chosen on $12$ tuning instances per class, so subject to its own winner's curse), and the action chosen without the size of the tested instances gains beyond the margin at all four sizes. The importance of the size found with the scheduling benchmark (Section~\ref{sec:granularity}) is therefore a property of that host, whose best action changes with $n$, not of the protocol; what carries over is that a static action chosen offline, at whatever granularity, beats every online learner.

% Auto-generated by analysis/make_table_journal.py
\begin{table}[t]
\caption{A second domain: the iterated local search for the Euclidean TSP. Gain over random choice (percentage points of the tour length, positive = better; $50$ test instances per class, $T=1000$, three seeds) of the validated action (val), the standard Q-Learning (ql), Q-Learning with the actions permuted in each run (ql\_perm), UCB1, the policy gradient (pg), the ALNS roulette (alns), Thompson sampling (ts) and the optimistic Q-Learning (qo). T5: five perturbation sizes; T15: three perturbations $\times$ five sizes. Bold: gain confirmed (Holm over the eight classes of each method and space) and not equivalent to zero at the margin $\pm0.05$. No online learner beats the validated action beyond the margin in any of the $96$ learner--cell pairs.}
\label{tab:tsp}
\centering
{\scriptsize\setlength{\tabcolsep}{4pt}
\begin{tabular}{@{}llrrrrrrrr@{}}
\toprule
Space & Class & val & ql & ql\_perm & ucb & pg & alns & ts & qo \\
\midrule
T5 & clustered $n=200$ & \textbf{0.12} & $-$0.15 & $-$0.05 & 0.03 & $-$0.02 & 0.04 & 0.05 & 0.02 \\
T5 & clustered $n=500$ & \textbf{0.59} & $-$0.46 & 0.00 & 0.17 & 0.22 & 0.04 & 0.27 & 0.23 \\
T5 & clustered $n=1000$ & \textbf{0.51} & $-$0.63 & $-$0.17 & $-$0.12 & $-$0.08 & $-$0.14 & \textbf{0.24} & 0.00 \\
T5 & clustered $n=2000$ & \textbf{0.35} & $-$0.53 & $-$0.02 & 0.06 & 0.02 & $-$0.07 & \textbf{0.38} & \textbf{0.20} \\
T5 & uniform $n=200$ & \textbf{0.08} & $-$0.34 & $-$0.11 & $-$0.05 & $-$0.09 & $-$0.12 & 0.01 & $-$0.05 \\
T5 & uniform $n=500$ & \textbf{0.29} & $-$0.38 & $-$0.14 & $-$0.02 & $-$0.01 & $-$0.10 & 0.08 & 0.04 \\
T5 & uniform $n=1000$ & \textbf{0.49} & $-$0.40 & $-$0.07 & 0.05 & \textbf{0.06} & $-$0.05 & \textbf{0.25} & \textbf{0.11} \\
T5 & uniform $n=2000$ & \textbf{0.55} & $-$0.42 & 0.03 & \textbf{0.07} & 0.01 & $-$0.02 & \textbf{0.30} & \textbf{0.13} \\
T5 & mean & 0.37 & $-$0.41 & $-$0.07 & 0.02 & 0.01 & $-$0.05 & 0.20 & 0.08 \\
\midrule
T15 & clustered $n=200$ & \textbf{0.05} & $-$0.13 & $-$0.07 & $-$0.01 & $-$0.03 & $-$0.08 & $-$0.01 & 0.00 \\
T15 & clustered $n=500$ & \textbf{0.36} & $-$0.37 & $-$0.13 & 0.02 & 0.10 & $-$0.13 & 0.16 & 0.14 \\
T15 & clustered $n=1000$ & \textbf{0.77} & $-$0.53 & $-$0.16 & 0.07 & $-$0.01 & $-$0.20 & 0.11 & 0.05 \\
T15 & clustered $n=2000$ & \textbf{0.90} & $-$0.48 & $-$0.09 & 0.00 & $-$0.04 & $-$0.15 & \textbf{0.28} & \textbf{0.09} \\
T15 & uniform $n=200$ & \textbf{0.09} & $-$0.15 & $-$0.05 & 0.04 & 0.03 & 0.01 & 0.06 & 0.03 \\
T15 & uniform $n=500$ & \textbf{0.39} & $-$0.32 & $-$0.12 & 0.02 & 0.02 & $-$0.13 & 0.04 & 0.00 \\
T15 & uniform $n=1000$ & \textbf{0.67} & $-$0.29 & $-$0.10 & 0.02 & 0.06 & $-$0.13 & \textbf{0.18} & 0.08 \\
T15 & uniform $n=2000$ & \textbf{0.80} & $-$0.34 & 0.02 & 0.05 & 0.02 & $-$0.08 & \textbf{0.24} & 0.07 \\
T15 & mean & 0.50 & $-$0.33 & $-$0.09 & 0.03 & 0.02 & $-$0.11 & 0.13 & 0.06 \\
\bottomrule
\end{tabular}}
\end{table}


% Auto-generated by analysis/make_table_journal.py
\begin{table}[t]
\caption{The TSP host on $37$ published TSPLIB instances ($200\le n\le2103$, EUC\_2D; $T=1000$, three seeds): gain over random choice (percentage points of the tour length, positive = better; unit: the instance) and mean gap to the optimal tour length of the TSPLIB (\%). The static actions were chosen on the synthetic tuning instances of Section~\ref{sec:tsp}, not on TSPLIB instances. Bold: gain over random choice confirmed (Holm over the nine methods of each space) and not equivalent to zero at the margin $\pm0.05$. No online learner beats the static action of the nearest size beyond the margin.}
\label{tab:tsplib}
\centering
{\scriptsize\setlength{\tabcolsep}{5pt}
\begin{tabular}{@{}lrrrr@{}}
\toprule
 & \multicolumn{2}{c}{T5} & \multicolumn{2}{c}{T15} \\
\cmidrule(lr){2-3}\cmidrule(lr){4-5}
Method & gain & gap & gain & gap \\
\midrule
random choice & 0 & 1.93 & 0 & 1.79 \\
static, chosen on synthetic instances of the nearest size & \textbf{0.53} & 1.38 & \textbf{0.56} & 1.21 \\
static, one action chosen on synthetic instances & \textbf{0.53} & 1.38 & \textbf{0.68} & 1.08 \\
standard Q-Learning & $-$0.61 & 2.56 & $-$0.34 & 2.14 \\
Q-Learning with permuted actions & $-$0.15 & 2.08 & $-$0.19 & 1.98 \\
UCB1 & 0.04 & 1.88 & 0.05 & 1.73 \\
policy gradient & $-$0.04 & 1.97 & 0.03 & 1.76 \\
ALNS roulette & $-$0.02 & 1.95 & $-$0.11 & 1.90 \\
Thompson sampling & 0.12 & 1.80 & \textbf{0.24} & 1.54 \\
optimistic Q-Learning & 0.02 & 1.91 & 0.07 & 1.72 \\
\bottomrule
\end{tabular}}
\end{table}

